\documentclass[11pt]{article}

\usepackage[margin=1in]{geometry}
\usepackage{amsmath,amssymb,amsthm}
\usepackage{booktabs,longtable,array,multirow,tabularx}
\usepackage{hyperref}
\usepackage{xcolor}
\usepackage{enumitem}
\usepackage{listings}
\usepackage{graphicx}
\usepackage{microtype}

\hypersetup{
  colorlinks=true,
  linkcolor=blue!60!black,
  citecolor=blue!60!black,
  urlcolor=blue!60!black
}

\lstset{
  basicstyle=\ttfamily\scriptsize,
  breaklines=true,
  frame=single,
  columns=fullflexible
}

\newtheorem{definition}{Definition}
\newtheorem{lemma}{Lemma}
\newtheorem{proposition}{Proposition}
\newtheorem{claim}{Claim}
\newtheorem{remark}{Remark}

\title{Finite-Atlas Holonomy from Binary Self-Observation:\\
A Reproducible $S_3$ Consensus Event in Boundary Atlases}
\author{SIQ / RelGauge Research Program}
\date{Draft manuscript -- submission-fix revision, 2026-06-20}

\begin{document}
\maketitle

\begin{abstract}
We present a finite computational framework in which gauge structure is not assigned as a primitive symmetry group but is read from the consistency of finite boundary observations.  The primitive object is a finite transition system observed through a finite boundary atlas.  A chart atlas induces local descriptions of unresolved fibers; transition maps between overlapping charts produce loop holonomies.  Even very small finite systems can generate order-dependent reconciliation laws when multiple local chart disagreements overlap.  We report a reproducible binary ($q=2$) example in which information-preserving self-observation robustly generates finite $C_2$ chart holonomy and, at a critical birth window, two overlapping $C_2$ chart-cycle flips in the same namespace generate an exact nonabelian group $S_3$.  The nonabelian event is localized on a $17$-microstate support inside a $512$-state space.  It is not a stable equilibrium sector; a local pulse-train audit classifies it as a one-tick critical flash.  However, after freezing the transition table at the critical iteration, the event is algorithmically stable across the tested reconstruction family: $64/64$ reconstructions recover exact target and support-level $S_3$, with support Jaccard overlap $1.0$, and a parameter-scale audit over $48$ observation-method configurations also recovers exact target/support $S_3$ in every tested configuration.  Negative controls are central to the result: lateral certified $C_3/C_2$ interfaces remain product-like, paired-return and moving-carrier audits do not produce nonvacuous $S_3$, and endpoint/schedule-cut crossing-phase audits do not expose a deterministic three-phase carrier.  The strongest supported conclusion is therefore not that Standard Model gauge has been derived, nor that $S_3$ is generic, but that finite, information-preserving self-observation can produce a localized, structural, nonabelian reconciliation event without inserting a nonabelian group by hand.
\end{abstract}

\tableofcontents
\newpage

\section{Introduction}

Gauge theory is usually formulated by postulating a symmetry group and requiring invariance under its local action.  In the Standard Model, the gauge group is $SU(3)\times SU(2)\times U(1)$; in lattice gauge theory, group variables are placed on links; in fiber-bundle formulations, a structure group is specified and connections are then studied.  This paper studies the converse direction in a finite setting: can nontrivial gauge structure be generated by finite self-observation, without specifying a nonabelian group in advance?

The framework considered here begins with a simple structural premise: every observation is made through a finite boundary.  A finite boundary cannot resolve the entire microscopic state; it partitions the state space into visible classes and unresolved fibers.  Multiple partial boundary windows yield overlapping local descriptions of the same fiber.  Consistency of these local descriptions requires transition maps on overlaps.  Around chart cycles these maps may fail to compose trivially, producing holonomy.  This is the finite-atlas analogue of a gauge connection.

This situates the construction near, but not identical to, several established programs.  Wilson's lattice formulation and Kogut's review show how gauge degrees of freedom can be represented finitely on links once a gauge group is chosen \cite{wilson1974,kogut1979}.  Wheeler's ``it from bit'' program motivates the information-theoretic direction \cite{wheeler1990}.  Causal-set theory and Wolfram-style discrete rewriting models explore spacetime or dynamics from finite/discrete primitives \cite{sorkin2003,wolfram2020}.  The distinction here is that the finite group is not specified on links in advance; it is extracted post hoc from chart-overlap reconciliation in a finite observer atlas.

The central result of this manuscript is a reproducible binary example.  In a $q=2$, $n=9$ state space with $512$ microstates, an information-preserving self-observation iteration produces a critical transition table.  At iteration $8$ of one reproducible random-full-permutation candidate, a local namespace $(\mathrm{parent\ domain}=75,\mathrm{fiber}=7)$ contains two distinct order-two chart-cycle loop maps, $(0\;1)$ and $(0\;2)$, acting on a common three-label support.  Since two transpositions sharing exactly one element generate $S_3$, this is an exact nonabelian finite gauge closure.  The microstate support of the event has size $17$.

The paper makes three kinds of claims.

\begin{enumerate}[label=(\roman*)]
\item \textbf{Mathematical claims.}  These are ordinary finite group statements, such as the fact that two overlapping transpositions generate $S_3$.
\item \textbf{Computational claims.}  These are claims about the outputs of deterministic audit programs on specified inputs, seeds, and parameters.  They are reproducible from the accompanying code and artifacts.
\item \textbf{Interpretive claims.}  These connect the computational findings to the idea of matter as multi-observer consensus.  These are hypotheses and are deliberately stated more cautiously than the computational claims.
\end{enumerate}

The strongest supported statement is:
\begin{quote}
A finite binary self-observation pipeline with information-preserving effective dynamics produces a localized $17$-state structural $S_3$ consensus event, without inserting $S_3$ by hand.
\end{quote}

The paper does not claim to derive the Standard Model, particle masses, or continuous spacetime.  It does claim that a concrete nonabelian gauge-reconciliation event can arise from finite boundary consistency alone.

It is also useful to separate two uses of the word ``gauge.''  The computations demonstrate finite-atlas holonomy, including reproducible $C_2$ chart holonomy in information-preserving runs.  Whether these finite holonomies become physical gauge fields in the Yang--Mills, lattice-gauge, or bundle-theoretic sense is a further scaling and interpretation problem, not a conclusion of this paper.


\subsection*{Boxed finite-boundary axiom}
\fbox{\begin{minipage}{0.92\linewidth}
\textbf{Finite-boundary axiom.} Every observation is made through a finite boundary map
\(\pi:S\to\Sigma\), where \(|\Sigma|<|S|\) for nontrivial observations.  The inverse images
\(F_\sigma=\pi^{-1}(\sigma)\) are unresolved fibers.  A finite observer may use multiple
partial chart windows on these fibers, and descriptions on overlaps must be mutually
reconcilable by finite transition maps.  No external reservoir may absorb distinctions lost
by the boundary; thus information-preserving effective dynamics is the axiom-compatible
case, while representative collapse is an explicitly lossy control.
\end{minipage}}

The computational pipeline below is one finite realization of this axiom.  The manuscript's
main claim is constructive, not universal: it exhibits one reproducible finite system in which
this pipeline yields a localized structural \(S_3\) consensus event.  It does not prove that
every finite observer must generate \(S_3\), nor that the event is a particle or a Standard
Model gauge sector.

\section{Formal setup}

\subsection{Finite state systems}

Let
\[
S=\{0,1,\dots,q-1\}^n
\]
be a finite state space.  In the main result, $q=2$ and $n=9$, so $|S|=512$.  A microscopic dynamics is a map
\[
T:S\to S.
\]
The principal experiments use random full permutations as initial dynamics, so $T$ is bijective.  This choice is not a gauge input: it fixes a generic information-preserving initial transition table.

A finite observer does not access $s\in S$ directly.  It observes a boundary projection
\[
\pi:S\to \Sigma,
\]
where $\Sigma$ is a finite label set.  For $\sigma\in \Sigma$, the fiber
\[
F_\sigma=\pi^{-1}(\sigma)
\]
contains microscopic states unresolved by the observer.

\begin{definition}[Finite boundary atlas]
A finite boundary atlas $A$ consists of a finite set of domains, each carrying a partition of a subset of $S$, together with charts on selected fibers and overlap transition maps between chart labels where charts describe common fiber states.
\end{definition}

The implemented atlas is generated by recursive proliferation.  Starting from a simple boundary label, fibers are recursively refined according to predictive distinctions induced by the transition table.  The procedure is finite because the atlas capacity, maximum domains, maximum charts per fiber, and maximum chart coordinate/support sizes are bounded.

\subsection{Charts, overlap maps, and holonomy}

Let $F$ be a fiber, and let $c_i:F\to L_i$ be a chart assigning finite chart labels to states in $F$.  On a region where two charts describe the same states, the induced transition map is a partial map
\[
g_{ij}:L_i\dashrightarrow L_j.
\]
A chart cycle is a sequence
\[
c_{i_0}\to c_{i_1}\to\cdots\to c_{i_m}=c_{i_0}.
\]
The composed loop map is
\[
H=g_{i_{m-1}i_m}\circ\cdots\circ g_{i_0 i_1}.
\]
When $H$ is a nontrivial permutation of chart labels, the atlas has finite holonomy.  An order-two loop map is a $C_2$ holonomy.

\begin{definition}[Same-namespace composition]
Two loop maps are composable for gauge closure only if they act on the same atlas namespace: the same candidate, iteration window, parent domain, and fiber label, or a certified equivalent support.  This requirement avoids conflating labels from different observers or different fibers.
\end{definition}

\subsection{Information-preserving self-observation}

Given an atlas $A(T)$ built from a transition table $T$, the pipeline constructs an effective transition table
\[
T_{\mathrm{eff}}=\Phi(T)
\]
by reading how $T$ acts on atlas labels and lifting the finite effective dynamics back to sampled microstates.  The next self-observation step uses
\[
T\leftarrow T_{\mathrm{eff}}.
\]

Several lift names are implemented: \texttt{bijective}, \texttt{permutation}, \texttt{conservative}, \texttt{information\_conserving}, and \texttt{representative}.  The first four preserve the relevant information in the reference regime and produce identical gauge-generation statistics; the representative lift collapses fibers to representatives and destroys gauge generation in the reference test.  Thus the empirical split is between information-preserving and information-destroying lifts, not between arbitrary tunable choices.

\section{Finite group facts}

The nonabelian result rests on an elementary theorem.

\begin{lemma}[Overlapping transpositions generate $S_3$]
Let $X=\{a,b,c\}$ with distinct elements.  Let $\tau_1=(a\;b)$ and $\tau_2=(a\;c)$.  Then
\[
\langle \tau_1,\tau_2\rangle\cong S_3.
\]
\end{lemma}

\begin{proof}
Both generators are transpositions.  Their product is a 3-cycle:
\[
(a\;b)(a\;c)=(a\;c\;b)
\]
up to composition convention.  The subgroup generated by a transposition and a 3-cycle on the same three-element set contains all six permutations of $X$.  Hence it is $S_3$.
\end{proof}

\begin{proposition}[Finite-atlas $S_3$ certificate]
If two same-namespace chart-cycle loop maps are transpositions sharing exactly one chart label, then the finite chart-cycle closure contains exact $S_3$.
\end{proposition}

\begin{proof}
Restrict to the union of the three labels involved.  Extend each transposition by the identity on labels outside its two-point support.  The preceding lemma applies.
\end{proof}

This proposition is the algebraic certificate used in the spatial chart-cycle closure and transposition-diversity audits.

\section{Computational methods and reproducibility}

\subsection{State space and reference candidate}

The reference nonabelian event occurs in a binary state space with $n=9$ vertices.  The target candidate is
\begin{center}
\texttt{random\_full\_permutation|inst=28|full\_atlas|cap=32|seed=3221741}.
\end{center}
The event occurs at atlas iteration $8$, in parent domain $75$, fiber label $7$.

\subsection{Deterministic seeding}

A reproducible implementation must not use Python's process-randomized built-in \texttt{hash()} in seed construction.  The patched iteration code uses deterministic SHA1 hashing of profile strings.  The relevant source file includes \texttt{hashlib} and the deterministic seed path.  The repository tests include a deterministic-hash check and grep for process-randomized seed paths.

\subsection{Frozen transition table}

For multi-observer tests, the transition table at the critical iteration is saved as a \texttt{.npy} array.  Downstream consensus audits load this frozen transition directly rather than replaying the entire iteration history.  This removes ambiguity due to earlier non-deterministic hash behavior and makes observer reconstructions act on the same boundary configuration.

\subsection{Audit modules}

The paper relies on the following audit families.

\begin{itemize}
\item \textbf{Iterated fiber-atlas dynamics:} generates atlas sequences and tracks $C_2$ holonomy.
\item \textbf{Spatial chart-cycle closure and transposition diversity:} composes same-namespace $C_2$ loop maps and classifies generated finite groups.
\item \textbf{Critical birth-window and pulse-train audits:} locate local $C_2$ birth/rebound windows and classify flat, $C_2$, and $S_3$ pulses.
\item \textbf{Critical-event microscope:} recomputes the $S_3$ group, orbit size, and matched $C_2$ controls.
\item \textbf{Multi-observer consensus:} rebuilds atlas observations from a frozen transition table and measures exact-namespace and support-level consensus.
\item \textbf{Observation-method scale test:} varies observer profile, atlas capacity, and entropy threshold while measuring whether the $S_3$ support remains visible.
\item \textbf{Negative-control audits:} certified inter-scale interfaces, paired returns, moving carriers, and crossing-phase tests.
\end{itemize}

\subsection{Result provenance}

Table~\ref{tab:provenance} lists the principal result artifacts used in this paper.  Filenames in the package and archive use the longer machine-readable names indicated in the repository manifest.

\begin{table}[p]
\centering
\footnotesize
\caption{Principal computational artifacts.}\label{tab:provenance}
\begin{tabularx}{\textwidth}{>{\raggedright\arraybackslash}p{0.34\textwidth}X}
\toprule
Artifact family & Role \\
\midrule
Witness lift-mode summaries & Reference 29-instance $q=2$, $v=9$, full-atlas, capacity-32 lift-mode comparison. \\
Broad $v=9$ iterated atlas run & Gauge-generation and spatial-closure source run. \\
Broad $v=10$ iterated atlas run & Scaling/control run showing strong $C_2$ generation but no $S_3$ closure. \\
Spatial closure and transposition-diversity audits & Same-namespace composition of $C_2$ chart loops; detects one exact $S_3$ event in $v=9$. \\
Critical birth-window audit & Birth-window provenance; detects one $S_3$ critical window among $C_2$ birth-window namespaces. \\
Critical-event microscope & Group/orbit certificate and matched $C_2$ controls. \\
Critical pulse train & Temporal localization of the $S_3$ event as a one-tick flash. \\
Frozen multi-observer consensus & 64-run exact/support $S_3$ consensus on the frozen transition. \\
Observation-method scale test & 48-configuration observation-method scale test. \\
\bottomrule
\end{tabularx}
\end{table}

\section{Results I: information preservation and $C_2$ generation}

\subsection{Lift-mode split}

The reference lift-mode experiment used $q=2$, $v=9$, random full permutations, $29$ instances, full-atlas profile, capacity $32$, $12$ atlas iterations, and seed control sufficient to reproduce the reference witness.  The information-preserving modes \texttt{bijective}, \texttt{permutation}, \texttt{conservative}, and \texttt{information\_conserving} produced the same summarized behavior: $C_2$ generation fraction $20/29=0.689655\ldots$, maximum chart $C_2$ holonomy $116$, and identical domain and chart-cycle counts in the reported run.  The \texttt{representative} lift produced no $C_2$ gauge generation and collapsed to flat/fixed behavior in the reference comparison.

\begin{table}[h]
\centering
\caption{Reference lift-mode result.  The representative lift is information-destroying and does not generate gauge in this test.}
\begin{tabular}{lrrrr}
\toprule
Lift mode & $C_2$ generation & Final $C_2$ & Max $C_2$ holonomy & Interpretation \\
\midrule
bijective & $0.689655$ & $0.0$ & $116$ & information-preserving \\
permutation & $0.689655$ & $0.0$ & $116$ & information-preserving \\
conservative & $0.689655$ & $0.0$ & $116$ & information-preserving in this regime \\
information\_conserving & $0.689655$ & $0.0$ & $116$ & information-preserving \\
representative & $0.0$ & $0.0$ & $0$ & information-destroying \\
\bottomrule
\end{tabular}
\label{tab:lift}
\end{table}

This result is important because it reduces an apparent implementation choice to an empirical equivalence class.  In the tested regime, gauge generation is not tied to one named lift.  It appears for the information-preserving class and vanishes for the representative collapse.

\subsection{Broad $C_2$ generation}

In the broad $v=9$ run, $C_2$ generation occurs in $0.41944$ of candidate histories, with final persistence $0.03148$ and maximum chart $C_2$ holonomy $704$.  In the broad $v=10$ run, $C_2$ generation rises to $0.58194$, final persistence to $0.16528$, and maximum chart $C_2$ holonomy to $800$.  Thus $C_2$ generation is robust and increases in the tested $v=10$ run.

However, $C_2$ abundance is not the same as nonabelian closure.  The $v=10$ spatial closure audit finds $4812$ same-namespace $C_2$ groups, all single or duplicate flips, with no shared-label pair, no component of size $3$, and no $S_3$.  This is a key negative control: more $C_2$ does not automatically produce $S_3$.

\section{Results II: exact $S_3$ from spatial chart-cycle closure}

\subsection{The spatial closure event}

The $v=9$ spatial chart-cycle closure audit finds $3838$ same-namespace $C_2$-bearing groups.  Of these, $3837$ are single or duplicate $C_2$ flips, and exactly one has two overlapping transpositions.  The exact $S_3$ fraction is therefore
\[
\frac{1}{3838}\approx 2.6055\times 10^{-4}.
\]
The maximum generated group order is $6$, and the maximum transposition component size is $3$.

The witness occurs at
\begin{align*}
\mathrm{candidate} &= \texttt{random\_full\_permutation|inst=28|full\_atlas|cap=32|seed=3221741},\\
\mathrm{iteration} &= 8,\\
\mathrm{parent\ domain} &= 75,\\
\mathrm{fiber\ label} &= 7.
\end{align*}

The nine nontrivial $C_2$ chart cycles in this namespace reduce to two distinct loop maps:
\[
(0\;2),\qquad (0\;1).
\]
Thus, by Proposition 1, the local chart-cycle closure is exact $S_3$.

\subsection{Route redundancy}

The event is not supported by a single accidental loop pair.  The $(0\;2)$ map appears along three chart routes,
\[
1\to 5\to 8\to 7,
\quad
1\to 5\to 9\to 7,
\quad
1\to 5\to 10\to 7,
\]
while the $(0\;1)$ map appears along six routes,
\[
5\to 8\to 7\to 13,
\quad
5\to 9\to 7\to 13,
\quad
5\to 10\to 7\to 13,
\]
\[
7\to 8\to 12\to 13,
\quad
7\to 9\to 12\to 13,
\quad
7\to 10\to 12\to 13.
\]
This route redundancy strengthens the interpretation that the event is a structural chart-consensus feature of the frozen transition, not a single-cycle parsing accident.

\section{Results III: critical birth-window provenance}

The critical-birth-window audit restricts to namespaces in the first $C_2$ birth window, with generated-from-flat dynamics.  It finds
\[
1164\quad C_2\text{-single/duplicate namespaces},
\]
and
\[
1\quad S_3\text{ critical-birth namespace}.
\]
The unique $S_3$ row has two distinct $C_2$ maps, a shared-label transposition pair, maximum transposition component size $3$, generated group order $6$, and a compression-rebound score approximately $359.50$.

Matched controls in the critical-event microscope have group order $2$ and orbit size $2$, whereas the target has group order $6$ and orbit size $3$.  The mean compression-rebound score of controls is approximately $25.05$, far below the target's score.

\begin{table}[h]
\centering
\caption{Critical-event microscope: target versus matched $C_2$ controls.}
\begin{tabular}{lrrrr}
\toprule
Class & Count & Group order & Largest orbit & Mean rebound score \\
\midrule
Target $S_3$ event & 1 & 6 & 3 & 359.50 \\
Matched $C_2$ controls & 20 & 2 & 2 & 25.05 \\
\bottomrule
\end{tabular}
\label{tab:microscope}
\end{table}

The matched-control count in Table~\ref{tab:microscope} is taken from the current verified microscope summary: $n_{\mathrm{controls}}=20$, all with group order $2$, and mean control compression-rebound score $25.05$.  An earlier exploratory microscope run used $10$ controls and is superseded by this verified artifact.

The pulse-train audit shows that this event is temporally local.  In the local namespace, the sequence is
\[
7:\mathrm{flat}\quad\to\quad 8:S_3\mathrm{\ flash}\quad\to\quad 9:\mathrm{flat}.
\]
The $S_3$ duration is one atlas tick, with collapse mode \texttt{to\_flat}.  Thus the event is not a stable equilibrium sector but a local critical consensus flash.

\section{Results IV: the 17-microstate support}

The target support contains $17$ microstate indices:
\[
\{29,30,44,47,61,62,156,159,188,191,221,222,253,254,256,284,287\}.
\]
Table~\ref{tab:support} lists their binary representations and transition images under the frozen transition table at iteration $8$.

\begin{longtable}{rllr}
\caption{The $17$ microstates of the structural $S_3$ support.}\label{tab:support}\\
\toprule
Index & Binary state & Image index & Hamming weight \\
\midrule
\endfirsthead
\toprule
Index & Binary state & Image index & Hamming weight \\
\midrule
\endhead
29 & 000011101 & 33 & 4 \\
30 & 000011110 & 34 & 4 \\
44 & 000101100 & 48 & 3 \\
47 & 000101111 & 51 & 5 \\
61 & 000111101 & 65 & 5 \\
62 & 000111110 & 66 & 5 \\
156 & 010011100 & 160 & 4 \\
159 & 010011111 & 163 & 6 \\
188 & 010111100 & 96 & 5 \\
191 & 010111111 & 99 & 7 \\
221 & 011011101 & 225 & 6 \\
222 & 011011110 & 226 & 6 \\
253 & 011111101 & 129 & 7 \\
254 & 011111110 & 130 & 7 \\
256 & 100000000 & 329 & 1 \\
284 & 100011100 & 288 & 4 \\
287 & 100011111 & 291 & 6 \\
\bottomrule
\end{longtable}

The support has a visible $16+1$ structure: eight paired states plus a singleton $256=100000000$.  The support is not dynamically closed: its image under the frozen transition lies outside the support.  Therefore the object is best described as a localized event support, not a persistent orbit of the microscopic transition.

\section{Results V: multi-observer and observation-method consensus}

\subsection{Frozen-transition multi-observer consensus}

The frozen-transition multi-observer audit loads the saved transition table at iteration $8$ and rebuilds the atlas $64$ times with independent observer seeds.  In all $64$ reconstructions, the target namespace is observed and exact $S_3$ is recovered.  Support-overlap matching also recovers exact $S_3$ with Jaccard overlap $1.0$.

The summary reports:
\begin{itemize}
\item target exact $S_3$ consensus fraction: $1.0$;
\item support exact $S_3$ consensus fraction: $1.0$;
\item target support size: $17$;
\item maximum target/support generated group order: $6$;
\item maximum target/support transposition component size: $3$;
\item observer-dependent candidate: false.
\end{itemize}

This is the strongest evidence that the $S_3$ event is not merely one observer's chart artifact.  Under the tested reconstruction family, the frozen transition forces the observer to recover the same nonabelian support.

\subsection{Observation-method scale test}

The observation-method scale test varies observer profile, atlas capacity, and chart entropy threshold over $48$ tested configurations.  Each configuration uses $16$ observer reconstructions.  All $48$ configurations recover exact target and support-level $S_3$, with support Jaccard $1.0$, support generated group order $6$, and maximum transposition component size $3$.

The tested grid includes observer profiles \texttt{full\_atlas}, \texttt{fiber\_preserving}, and \texttt{boundary\_only}; capacities $16,32,64,128$; and minimum chart entropy thresholds $0,0.025,0.05,0.1$, with fixed chart-coordinate settings in the reported small grid.  The result should therefore be stated as invariance across this tested observation-method family, not across all possible observers.

\section{Negative results and controls}

Negative results are essential because they locate the actual mechanism.

\subsection{Certified $C_3/C_2$ lateral interface is product-like}

A certified-sector interface audit constructs a certified $C_3$ event sector and a certified $C_2$ frame sector.  The certified interface exists in the twist/hybrid sectors, but it remains product-like.  The summary reports certified-interface fraction $0.4$, product-like fraction $0.4$, compressed fraction $0$, coupled-candidate fraction $0$, product-subtracted candidate fraction $0$, maximum period $6$, and maximum bijective domain $6$.  Thus the side-by-side interface realizes the product baseline $C_3\times C_2\cong C_6$, not a semidirect $S_3$ extension.

\subsection{Paired return and moving carriers do not produce exact $S_3$}

Paired-return audits construct a shared carrier $(\sigma,f)$ and test observed return maps.  They find either empty carriers, size-two cores, weak feedback, or vacuous relations, but no order-three return, no noncommuting return, and no strict $S_3$.  Moving-carrier audits find transport arrows and weak $C_2$ naturality failure, but again no order-three moving return, no noncommuting cycle, and no strict $S_3$.

\subsection{Endpoint and schedule-cut crossing phases are negative}

Endpoint and schedule-cut crossing-phase audits test whether a hidden pre/mid/post phase cycle yields $S_3$.  Both are negative.  Phase rows exist and $C_2$ flips can be present, but no three-phase triplet, no aggregate order-three map, and no aggregate $S_3$ appear.  This rules down the hypothesis that the nonabelian structure arises as a deterministic internal crossing phase.

\subsection{$v=10$ scaling control}

The broad $v=10$ run increases $C_2$ generation and final persistence, but its spatial closure remains a pure $C_2$ baseline: $4812$ same-namespace $C_2$ groups, all single/duplicate flips, no shared-label pair, no component of size $3$, and no exact $S_3$.  Thus the $S_3$ event is not a monotone consequence of more vertices or more $C_2$ abundance.  It requires transposition diversity in a local namespace.

\section{Interpretation: gauge, consensus, and event-like matter}

The results support a distinction between three layers:

\begin{enumerate}
\item \textbf{Gauge baseline:} robust $C_2$ chart holonomy generated by information-preserving finite self-observation.
\item \textbf{Nonabelian reconciliation event:} a local critical window where two incompatible $C_2$ chart disagreements overlap, forcing exact $S_3$.
\item \textbf{Matter-like consensus candidate:} a nontrivial support that independent observers recover as the same structure under support-level matching and observation-method variation.
\end{enumerate}

The verified $17$-state event satisfies the first two layers and passes the tested multi-observer consensus tests.  It is therefore reasonable to call it a localized structural $S_3$ consensus event.  It is not yet a particle in the usual physical sense because it has not been shown to propagate, recur, carry a conserved charge, or exhibit perturbation resistance.  It is event-like rather than object-like.

This suggests the following operational definitions.

\begin{definition}[Gauge disagreement]
Gauge disagreement is the groupoid generated by transition maps between distinct chart descriptions of the same boundary support.
\end{definition}

\begin{definition}[Matter-like consensus]
A matter-like consensus support is a localized microstate support whose nontrivial reconciliation group is recovered by independent observer reconstructions over a specified observation-method class.
\end{definition}

Under these definitions, the reference $17$-state support is the first matter-like consensus candidate in the framework.


\subsection*{Statistical status of the rare event}
The spatial closure rate is deliberately reported as rare: one exact \(S_3\) namespace among
3838 \(C_2\)-bearing namespaces in the v9 chart-cycle closure audit, and one exact critical
birth-window event among 1165 generated birth-window namespaces.  The claim in this paper is
therefore not a prevalence claim.  Frequency alone is not used as evidence.  The evidence is
that the rare event has an exact finite-group certificate, a localized 17-state support, a
matched-control contrast, 64/64 frozen-transition observer recovery, and 48/48 recovery in the
reported observation-method grid.  A future submission should add a constrained label-shuffle
or loop-map-shuffle null distribution for the exact 1/3838 count; the present draft treats this
as an explicit limitation rather than a completed significance test.


\subsection{Aggregate certified-sector near-miss is not positive evidence}
A separate certified-sector algebra audit reported a small aggregate ``S3-like'' C3/C2 interface signature without exact group proof.  This is not used as evidence for the main result.  The corresponding certified interface remains product-like: certified C3 and C2 sectors coexist on a six-state product baseline, with no compressed interface, no product-subtracted candidate, and no exact semidirect closure.  This near-miss is useful as a negative control because it shows that the present positive S3 event is not a generic artifact of any six-state C3/C2 product table.

\subsection{Paired-return verdict correction}
Early paired-return verdict text overstated a subcritical feedback diagnostic.  The corrected interpretation is that generated candidates can show size-two coordinate-feedback diagnostics, but no valid size-six joint return carrier, no order-three return, no noncommuting return, and no nonvacuous S3 relation.  These paired-return rows are therefore negative controls, not supporting evidence.

\section{Limitations}

The present result is strong but limited.

\begin{enumerate}
\item \textbf{Existence, not prevalence.}  The verified $S_3$ event is unique in the reported $v=9$ birth-window audit.  It is structurally verified but rare.
\item \textbf{Finite observation-method family.}  The scale test covers $48$ configurations, not all possible observers.  The result should be stated relative to that tested family.
\item \textbf{No Standard Model identification.}  $S_3$ is not $SU(2)$ or $SU(3)$.  The result demonstrates finite nonabelian gauge reconciliation, not Standard Model recovery.
\item \textbf{No propagation yet.}  The $17$-state support maps out of itself and the local $S_3$ pulse lasts one tick.  A memory-transfer or perturbation-stability audit is required before claiming particle-like dynamics.
\item \textbf{Implementation dependence.}  Although the package provides reproducibility tests, any computational result depends on the correctness of the implementation.  The paper mitigates this by exposing frozen transition tables, support hashes, synthetic tests, and simple finite-group certificates.
\end{enumerate}

\section{Reproducibility}

\subsection*{Repository scope}
The accompanying repository is a real computational package for the central witness and consensus tests.  It includes the actual \texttt{relgauge} modules needed for the core positive-result chain, the frozen transition table \texttt{data/transition\_iter8\_inst28\_seed3221741.npy}, the 17-state reference support metadata, shell scripts, tests, and the paper source.  It does not rely on JSON-summary wrapper scripts for the central frozen-transition checks.

The repository is intentionally scoped.  The script \texttt{reproduce.sh} runs the \emph{core positive-result reproduction chain}: witness generation, regenerated-transition comparison, critical birth-window detection, frozen-transition multi-observer consensus, observation-method scale, and lift-mode comparison.  It does not rerun every exploratory or negative-control audit described in the manuscript.  The broader negative controls--certified $C_3/C_2$ interfaces, paired returns, moving carriers, endpoint and schedule-cut crossing phases, broad $v=10$ scaling controls, and other exploratory modules--are reported from archived run artifacts and are not all included in the minimal core package.

\subsection*{Package commands}
The package uses Bash-style line continuation.  A reviewer can run:

\begin{lstlisting}[language=bash]
pip install -r requirements.txt
bash test_imports.sh
bash reproduce_quick.sh
\end{lstlisting}

The quick script recomputes the frozen-transition multi-observer and observation-method scale tests from the bundled transition table.  It should reproduce the target/support-level $S_3$ consensus without reading precomputed summary JSON files.

For the longer core chain:

\begin{lstlisting}[language=bash]
bash reproduce.sh
\end{lstlisting}

This can take tens of minutes to more than an hour depending on CPU and disk speed because it regenerates witness CSVs and chart-cycle rows.  The package also includes \texttt{requirements-lock.txt} and \texttt{environment.yml} recording the validated Python and dependency versions.

\section{Conclusion}

This paper demonstrates a concrete finite route from binary self-observation to nonabelian gauge reconciliation.  Information-preserving finite atlas iteration robustly generates $C_2$ chart holonomy.  In a critical local birth/rebound window, two overlapping same-namespace $C_2$ flips appear on a $17$-microstate support.  Their closure is exact $S_3$.  Freezing the transition at that moment and reconstructing the atlas independently yields exact target and support-level $S_3$ in all tested observer reconstructions, and the result survives all tested observation-method configurations.

The result does not derive the Standard Model.  It does not show that $S_3$ is generic.  It does show that a finite binary boundary-observation pipeline can produce localized, structural, nonabelian gauge consensus without inserting a nonabelian group by hand.  The next decisive tests are support-flux, transition-neighborhood consensus at iterations $7,8,9$, and perturbation stability of the $17$-state support.


\appendix
\section{Core algorithms in pseudocode}\label{app:pseudocode}

This appendix gives implementation-level pseudocode for the operations used in the reference audits.  The goal is not to replace the released code, but to make the mathematical construction explicit enough that a reader can reconstruct the finite objects independently.

\subsection{Iterated finite self-observation}
\begin{lstlisting}
Input: finite state list S, transition table T:S->S, alphabet q,
       atlas capacity K, iteration count M, observer parameters theta.
Output: atlas summaries A_t, effective transitions T_t.

T_0 := T
for t = 0,...,M:
    A_t := BUILD_ATLAS(S, T_t, theta)
    L_t := PROFILE_LABELS(A_t, T_t, profile)
    B_t := COMPRESS_LABELS_TO_CAPACITY(L_t, K)
    R_t := TEMPORAL_RELATION(B_t, T_t)
    T_{t+1} := INFORMATION_PRESERVING_LIFT(B_t, R_t, S, T_t)
    record domains, atlas maps, chart cycles, C2/S3 holonomies
end for
\end{lstlisting}

The reference experiment uses a deterministic seed schedule.  Python's process-randomized \texttt{hash()} is not used; profile-dependent seeds are computed by a stable SHA1 hash.  When \texttt{--save-transition-at t} is enabled, the array representing $T_t$ is written to disk before the atlas pass at iteration $t$.

\subsection{Atlas generation and recursive proliferation}
\begin{lstlisting}
BUILD_ATLAS(S, T, theta):
    D_0 := { root domain with initial boundary label pi_0(s) }
    D := D_0
    for depth = 1,...,max_depth:
        new_domains := []
        for domain d in D at previous depth:
            for each live fiber F of d with |F| >= min_fiber_size:
                compute predictive signatures of states in F under T over horizon h
                split F by signatures satisfying entropy/live-class filters
                keep at most max_domains_per_depth highest-scoring child domains
                append kept children to new_domains
        D := D union new_domains
    return domains D plus chart data from BUILD_CHARTS(D,T,theta)
\end{lstlisting}

A domain is therefore not a spatial region in a pre-existing space.  It is a finite predictive partition induced by the boundary projection and the transition table.  A fiber is a set of microstates sharing the same domain label.

\subsection{Chart construction on one fiber}
\begin{lstlisting}
BUILD_CHARTS_FOR_FIBER(domain d, fiber F, state table S, transition T):
    charts := []
    for coordinate subset C with |C| <= max_chart_coords:
        label_C(s) := tuple of selected boundary/support coordinates of s
        if label_C partitions F into at least min_chart_classes
           and entropy(label_C on F) >= min_chart_entropy
           and every retained class has enough support states:
              add chart c=(C,label_C,support=F)
    keep at most max_charts_per_fiber charts by entropy/support score
    return charts
\end{lstlisting}

The concrete implementation uses deterministic enumeration or seeded deterministic selection of coordinate subsets subject to the same filters.  Changing these admissible chart parameters is what the observation-method scale audit tests.

\subsection{Overlap maps and chart-cycle holonomy}
\begin{lstlisting}
OVERLAP_MAP(chart c_i, chart c_j):
    shared := support(c_i) intersection support(c_j)
    if |shared| < min_overlap_states: reject
    for each source label a of c_i on shared:
        count target labels b = c_j(s) over states s with c_i(s)=a
        if one target b has deterministic majority passing threshold:
            map a -> b
        else:
            mark source label nondeterministic
    return partial map g_ij on deterministic source labels

CHART_CYCLE_HOLONOMY(charts, overlap maps):
    build graph whose vertices are charts and edges are valid overlap maps
    enumerate simple cycles up to max_cycle_len
    for cycle c0->c1->...->ck=c0:
        H := g_{c_{k-1},c_k} o ... o g_{c_0,c_1}
        core := largest subset of labels closed under H and H^{-1}
        if H restricted to core is a permutation:
            order := permutation_order(H|core)
            record C2 if order=2; record loop map string
\end{lstlisting}

The same-namespace rule is applied after this step: two loop maps are composed for group closure only when they belong to the same candidate, atlas iteration/window, parent domain, and fiber label, or when a later support-overlap consensus test certifies that two observers are describing the same microstate support.

\subsection{Transposition diversity and exact $S_3$ closure}
\begin{lstlisting}
TRANSPOSITION_DIVERSITY(namespace N):
    maps := distinct order-2 loop maps recorded in N
    E := set of single-transposition edges {a,b} induced by maps
    build graph G_N with label vertices and transposition edges E
    component_size := size of largest connected component of G_N
    generated_group := closure of maps as finite permutations on union core
    exact_S3 := (|generated_group|=6 and group is nonabelian
                 and component_size=3)
    return distinct map count, component_size, exact_S3
\end{lstlisting}

For the reference event, $E=\{\{0,1\},\{0,2\}\}$, so the generated group is exactly $S_3$.

\subsection{Critical birth-window provenance}
\begin{lstlisting}
CRITICAL_BIRTH_WINDOW(namespace N at iteration t):
    first_window := |t - first_C2_iteration(candidate)| <= 1
    c2_birth := (C2_count(t) > 0 and C2_count(t-1) = 0)
    signature_delta := full_signature_classes(t) - full_signature_classes(t-1)
    fiber_entropy_delta := fiber_entropy(t) - fiber_entropy(t-1)
    valid_cycle_drop := max(0, valid_cycles(t-1)-valid_cycles(t)) / max(1, valid_cycles(t-1))
    rebound_score := max(0,signature_delta)
                     + 50*max(0,fiber_entropy_delta)
                     + 100*valid_cycle_drop
                     + (25 if c2_birth else 0)
    priority := 5*transposition_diversity_score + 0.01*rebound_score
                + domain-depth/live-fiber/predictive-class terms
    return provenance fields and exact_S3 classification
\end{lstlisting}

The rebound score is a heuristic provenance statistic, not a fitted physical observable.  It is used to compare the verified $S_3$ event to matched $C_2$ controls and to prioritize near misses.

\subsection{Multi-observer and support-overlap consensus}
\begin{lstlisting}
MULTI_OBSERVER_CONSENSUS(frozen transition T*, reference support R):
    for observer seed/configuration omega in tested family:
        A_omega := BUILD_ATLAS(S, T*, theta_omega)
        groups := all same-namespace transposition-diversity groups in A_omega
        exact_target := group at target parent/fiber is exact S3, if present
        for each group g in groups:
            J(g,R) := |support(g) intersection R| / |support(g) union R|
            overlap(g,R) := |support(g) intersection R| / |R|
        best := group maximizing J(g,R), subject to overlap thresholds
        record whether best is exact S3, C2-only, or flat
    report fractions over observer seeds/configurations
\end{lstlisting}

The term ``multi-observer'' in this paper therefore means algorithmic stability across a specified reconstruction family.  It is a finite model of observer consensus, not a proof of observer-independence in the physical or philosophical sense.

\section{Reproduction commands}

\subsection{Reference witness generation}

\begin{lstlisting}[language=bash]
python -m relgauge.iteratedfiberatlasdynamicsaudit 2 \
  --vertices 9 \
  --rule-modes random_full_permutation \
  --instances 29 \
  --atlas-capacities 32 \
  --profiles full_atlas \
  --atlas-iterations 12 \
  --proliferation-iterations 4 \
  --horizon 3 \
  --max-state-samples 512 \
  --save-transition-at 8 \
  --seed 2000006 \
  --out results/repro_witness_search.csv
\end{lstlisting}

\subsection{Critical birth-window audit}

\begin{lstlisting}[language=bash]
python -m relgauge.criticalc2birthwindowaudit \
  --chart-cycles-csv results/repro_witness_search_chart_cycles.csv \
  --iterated-csv results/repro_witness_search.csv \
  --domains-csv results/repro_witness_search_domains.csv \
  --require-generated --require-flat-start \
  --focus-first-c2-window 1 \
  --namespace-cols parent_domain_id,fiber_label \
  --out results/critical_birth.csv
\end{lstlisting}

\subsection{Multi-observer consensus}

\begin{lstlisting}[language=bash]
python -m relgauge.multiobserverconsensusaudit 2 \
  --vertices 9 \
  --iterated-csv results/repro_witness_search.csv \
  --frozen-transition-npy data/transition_iter8_inst28_seed3221741.npy \
  --target-rule-mode random_full_permutation \
  --target-instance 28 \
  --target-profile full_atlas \
  --target-atlas-capacity 32 \
  --target-seed 3221741 \
  --target-iteration 8 \
  --target-parent-domain 75 \
  --target-fiber-label 7 \
  --observer-runs 64 \
  --max-state-samples 512 \
  --out results/multi_observer_consensus.csv
\end{lstlisting}

\subsection{Observation-method scale test}

\begin{lstlisting}[language=bash]
python -m relgauge.observationmethodscaletest 2 \
  --vertices 9 \
  --iterated-csv results/repro_witness_search.csv \
  --frozen-transition-npy data/transition_iter8_inst28_seed3221741.npy \
  --target-rule-mode random_full_permutation \
  --target-instance 28 \
  --target-profile full_atlas \
  --target-atlas-capacity 32 \
  --target-seed 3221741 \
  --target-iteration 8 \
  --target-parent-domain 75 \
  --target-fiber-label 7 \
  --observer-profiles full_atlas,fiber_preserving,boundary_only \
  --observer-atlas-capacities 16,32,64,128 \
  --min-chart-entropy-list 0,0.025,0.05,0.1 \
  --observer-runs 16 \
  --max-state-samples 512 \
  --out results/observation_method_scale.csv
\end{lstlisting}

\section{Reference support hash}

The reference support is
\[
\{29,30,44,47,61,62,156,159,188,191,221,222,253,254,256,284,287\}.
\]
The support hash reported by the multi-observer audit is
\begin{center}
\texttt{ffb7d91caf4970bcfd60}.
\end{center}

\section{Recommended future audits}

\begin{enumerate}
\item \textbf{Transition-neighborhood consensus.}  Save and test transitions at iterations $7,8,9$ to verify that structural $S_3$ consensus is localized to iteration $8$.
\item \textbf{Support-flux memory audit.}  Track the image of the $17$-state support and test whether it overlaps later gauge-bearing supports.
\item \textbf{Perturbation stability.}  Modify transition rows inside/outside the support and measure the minimum perturbation that destroys $S_3$ consensus.
\item \textbf{Broader observation-method grid.}  Vary chart-coordinate limits, support-coordinate limits, and charts per fiber beyond the initial $48$ configurations.
\end{enumerate}

\begin{thebibliography}{9}

\bibitem{weyl1929}
H. Weyl, ``Elektron und Gravitation. I,'' \emph{Zeitschrift fuer Physik} \textbf{56}, 330--352 (1929).

\bibitem{yangmills1954}
C. N. Yang and R. L. Mills, ``Conservation of isotopic spin and isotopic gauge invariance,'' \emph{Physical Review} \textbf{96}, 191--195 (1954).

\bibitem{rovelli1996}
C. Rovelli, ``Relational quantum mechanics,'' \emph{International Journal of Theoretical Physics} \textbf{35}, 1637--1678 (1996).

\bibitem{isham1998}
C. J. Isham and J. Butterfield, ``A topos perspective on the Kochen-Specker theorem,'' \emph{International Journal of Theoretical Physics} \textbf{37}, 2669--2733 (1998).

\bibitem{hooft2016}
G. 't Hooft, \emph{The Cellular Automaton Interpretation of Quantum Mechanics}, Springer (2016).

\bibitem{wilson1974}
K. G. Wilson, ``Confinement of quarks,'' \emph{Physical Review D} \textbf{10}, 2445--2459 (1974).

\bibitem{kogut1979}
J. B. Kogut, ``An introduction to lattice gauge theory and spin systems,'' \emph{Reviews of Modern Physics} \textbf{51}, 659--713 (1979).

\bibitem{wheeler1990}
J. A. Wheeler, ``Information, physics, quantum: The search for links,'' in W. H. Zurek (ed.), \emph{Complexity, Entropy, and the Physics of Information}, Addison--Wesley, 3--28 (1990).

\bibitem{sorkin2003}
R. D. Sorkin, ``Causal sets: Discrete gravity,'' in A. Gomberoff and D. Marolf (eds.), \emph{Lectures on Quantum Gravity}, Springer, 305--327 (2005); arXiv:gr-qc/0309009 (2003).

\bibitem{wolfram2020}
S. Wolfram, ``A class of models with the potential to represent fundamental physics,'' \emph{Complex Systems} \textbf{29}, 107--536 (2020).

\bibitem{siqrepo}
SIQ / RelGauge Research Program, \emph{siq-framework reproducibility package}, included computational repository and verified result artifacts (2026).

\end{thebibliography}

\end{document}
